\subsection{COMMUTATIVELY CONVERGENT SERIES}

Let \((x_{n})\) be a summable sequence in \(G\) , and let \(s = \sum_{n\in \mathbb{N}}x_n\) be its sum. Then for each neighbourhood \(V\) of \(0\) , there exists \(J_0\in \mathfrak{F}(N)\) such that \(s_J\in s + V\) whenever \(J\in \mathfrak{F}(N)\) and \(J_0\subset J\) . Let \(m\) be the largest integer in \(J_0\) ; then if \(n\geqslant m\) we have \(s_n\in s + V\) , and therefore the series \((x_{n})\) is convergent and its sum is \(s\) . But the converse is false:

the sequence of terms of a convergent series can very well fail to be summable (see Chapter IV, § 7).

Moreover, the definition of a convergent series essentially involves the order structure of N. If the series \((x_{n})\) is convergent, and if \(\sigma\) is a permutation of N, then the series \((x_{\sigma(n)})\) is not necessarily convergent (cf.~Chapter IV, § 7, Exercise 15).

DEFINITION 2. A series defined by a sequence \((x_{n})\) is said to be commutatively convergent if, for each permutation \(\sigma\) of \(\mathbf{N}\) , the series defined by the sequence \((x_{\sigma(n)})\) is convergent.

PROPOSITION 9. The series defined by the sequence \((x_{n})\) is commutatively convergent if and only if the sequence \((x_{n})\) is summable; and then, for each permutation \(\sigma\) of \(\mathbf{N}\) , we have

\[
\sum_ {n = 0} ^ {\infty} x _ {\sigma} = \sum_ {n \in \mathbb {N}} x _ {n}.
\]

If the sequence is summable, then clearly the series is commutatively convergent. To prove the reverse implication we shall argue by reductio ad absurdum and suppose that the series \((x_{n})\) is commutatively convergent but that the sequence \((x_{n})\) is not summable. The image of the filter \(\Phi\) under the mapping \(\mathbf{H} \to s_{\mathbf{H}}\) therefore cannot be a Cauchy filter base in \(\mathbf{G}\) , otherwise this filter base would converge, since by hypothesis it has a cluster point (Chapter II, § 3, no. 2, Proposition 5, Corollary 2). Hence there is a neighbourhood \(\mathbf{V}\) of \(o\) in \(\mathbf{G}\) such that, for each finite subset \(\mathbf{J}\) of \(\mathbf{N}\) , there is a finite subset \(\mathbf{H}\) of \(\mathbf{N}\) which does not meet \(\mathbf{J}\) and is such that \(\sum_{n \in \mathbf{H}} x_{n} \notin V\) . Hence, by induction, we can define a partition of

N into finite subsets \(H_{k}\) ( \(k \in N\) ) such that \(\sum_{n \in H_{k}} x_{n} \notin V\) for an infinite number of indices k. Clearly there is a permutation \(\sigma\) of N such that for each k, the values of n for which \(\sigma(n) \in H_{k}\) are consecutive. If \(\sigma\) is such a permutation, then the series whose general term is \(x_{\sigma(n)}\) cannot be convergent, and we have a contradiction.

If the group G is written multiplicatively, the infinite product defined by a sequence \((x_{n})\) of points of G (or infinite product whose general term is \(x_{n}\) , or even product \(x_{n}\) if there is no possibility of confusion) is defined to be the pair formed by the sequence \((x_{n})\) and the sequence of partial products \(p_{n}=\prod_{k=0}^{n}x_{k}\) . The infinite product is said to converge if the sequence \((p_{n})\) converges, and the limit of this sequence is denoted by \(\prod_{n=0}^{\infty}x_{n}\) (or \(\prod_{n=0}^{\infty}x_{n}\) by abuse of notation). We leave to the reader the task of transcribing into multiplicative notation the properties of series which we have established.

